%%%%%%%%%%%%%%%%%%%%%%%%%%%%%%%%%%%%%%%%%%%%%%%%%%%%%%%%%%%%%%%%%%%%%%%%%%%%%%%%
%  Policy Evaluation Lab -- Lab 3: a case study, university scholarships (DSU)
%  Thursday 5 November 2026.
%  Source: MEF written exam, 31 March 2025, code STAT, envelope (busta) 2
%  (not drawn), question 2; published by the MEF on 4 June 2025 (ZIP of the
%  written-test papers on the competition page). Case text and questions
%  translated; Table 1 re-typed and checked: every column adds up to the
%  published total. Data constructed for the exercise (note kept).
%%%%%%%%%%%%%%%%%%%%%%%%%%%%%%%%%%%%%%%%%%%%%%%%%%%%%%%%%%%%%%%%%%%%%%%%%%%%%%%%

\documentclass[aspectratio=169,xcolor=table]{beamer}
\input{preamble-pel}

\title[Policy Evaluation Lab]{Policy Evaluation Lab}
\subtitle[Lab 3]{Lab 3: A case study, university scholarships}
\date[\pelYear]{November, 2026}

\begin{document}

\frame{\titlepage}

\begin{frame}{So, today we are going to evaluate a case study}
\be
  \item We look at a real case study
  \item Form groups of \textbf{2--3 people} at most
  \item Study what the question is asking, and look at the data
  \item Formulate some opinions, and the approach you would use
  \item Present your results
\ee
\end{frame}

\begin{frame}{Today's lab: a case study}
\small
Our country has a rather low share of people with a university degree and considerable inequality in access to tertiary education. Those who come from a family in economic difficulty, and whose parents have a low level of education, are less likely to enrol at university.

\vspace{0.15cm}
To encourage people to undertake university studies, the Italian State has introduced several measures. The main one is a \textbf{scholarship} for people who enrol in a degree course. To obtain it, there are both \textbf{merit} requirements and requirements linked to the \textbf{income} of the student's family. The management of the measures for the right to university study (\textit{diritto allo studio universitario}, DSU), meaning financial contributions and services for students, is delegated to the \textbf{Regions}. Recently this measure has been considerably strengthened.
\end{frame}

\begin{frame}{Table 1: scholarships by region}
\centering\footnotesize
\resizebox{!}{0.63\textheight}{%
\begin{tabular}[b]{l rrrrr}
\toprule
\textbf{Region} & \textbf{Scholarships} & \textbf{Eligible} & \textbf{Regular} & \textbf{\% of holders} & \textbf{Total spending} \\
 & \textbf{awarded} & \textbf{students} & \textbf{enrolled} & \textbf{away from home} & \textbf{(€ thousand)} \\
\midrule
Liguria & 3,264 & 3,264 & 25,108 & 43.8 & 8,205 \\
Lombardia & 25,861 & 25,933 & 243,972 & 44.4 & 78,952 \\
Piemonte & 15,843 & 15,843 & 95,440 & 46 & 48,219 \\
Valle d'Aosta & 131 & 131 & 814 & 35.9 & 293 \\
Emilia-Romagna & 22,607 & 22,607 & 142,182 & 39.1 & 76,124 \\
Friuli-Venezia Giulia & 5,286 & 5,286 & 24,248 & 37 & 18,612 \\
Trentino-Alto Adige & 2,902 & 2,902 & 14,510 & 50 & 8,309 \\
Veneto & 14,698 & 15,661 & 95,442 & 37 & 44,121 \\
Lazio & 21,959 & 21,959 & 246,730 & 66.6 & 73,636 \\
Marche & 5,227 & 5,227 & 35,318 & 81 & 21,613 \\
Toscana & 11,887 & 11,887 & 83,711 & 67.5 & 49,270 \\
Umbria & 4,506 & 4,506 & 23,227 & 55.2 & 15,446 \\
Abruzzo & 4,694 & 5,322 & 32,151 & 48.2 & 14,215 \\
Basilicata & 1,021 & 1,021 & 3,796 & 15.4 & 2,079 \\
Calabria & 10,282 & 10,798 & 27,714 & 66.9 & 36,226 \\
Campania & 27,543 & 27,543 & 120,803 & 6.6 & 111,646 \\
Molise & 902 & 1,081 & 5,214 & 16.5 & 3,116 \\
Puglia & 18,100 & 18,100 & 61,775 & 17.1 & 49,589 \\
Sardegna & 9,074 & 9,074 & 27,168 & 46.1 & 31,779 \\
Sicilia & 18,804 & 21,134 & 81,052 & 20.3 & 43,057 \\
\midrule
\textbf{Italy} & \textbf{224,591} & \textbf{229,279} & \textbf{1,390,375} & \textbf{40} & \textbf{734,507} \\
\bottomrule
\end{tabular}}
\pelTableNote{Students of degree courses in Italian universities, academic year 2021/2022; spending in calendar year 2021. ``Regular enrolled'': students enrolled within the normal duration of their course (\textit{iscritti regolari}). ``Away from home'': \textit{fuori sede}. Source: MEF written exam, 31 March 2025 (envelope 2, not drawn). The data were constructed for the exercise and do not necessarily correspond to real values.}
\end{frame}

\begin{frame}{The questions}
\begin{block}{}
Based on the information provided in Table 1, analyse the main evidence describing the process of awarding scholarships across the regions, indicating also which indicators it would be useful to compute, and how.
\end{block}
\vspace{0.2cm}
\begin{block}{}
In addition, illustrate which strategy could be adopted, and which other data would be needed, to evaluate the effectiveness of the scholarships for the right to university study in terms of improving the academic performance of enrolled students.
\end{block}
\end{frame}

\end{document}
